\section{Related Work}

The development of steady-state visual evoked potential (SSVEP)-based brain-computer interfaces (BCIs) has seen significant advancements, over the last decades, focusing on improving communication for individuals with restricted or none motor capabilities.

A notable contribution is the BCI Speller system developed by DTU \cite{dtubcispeller}, which integrates SSVEP responses with dictionary support for efficient text generation. This system utilizes a minimalistic setup with three electrodes and uses a two-stage selection model. Using a dictionary for word prediction, it significantly improves communication speed, achieving an average character production rate of 4.91 characters per minute (CPM). The DTU BCI Speller system demonstrates robust usability, even under non-ideal conditions, emphasizing user-friendly design and intuitive interfaces.

Complementing these efforts, recent advances in SSVEP detection have focused on deep learning methodologies. For example, a 2022 study published in \textit{Frontiers in Computational Neuroscience} introduced a "CNN-based approach to the detection of SSVEP using binaural EEG" \cite{cnnbased}. The mentioned method utilizes convolutional neural networks (CNNs) to extract relevant features from ear-centered EEG signals, offering a compact and portable alternative to traditional setups. By targeting the unique challenges of ear-EEG, this approach expands the applicability of SSVEP-based BCIs to scenarios requiring minimalistic and portable designs. The CNN-based model demonstrated superior performance in classifying SSVEP stimuli compared to conventional methods, showcasing the importance and potential of deep learning to improve detection accuracy and system adaptability.

Both of these works, collectively, demonstrate the rapid evolution of SSVEP-based BCIs driven by user-centric innovations and cutting-edge computational techniques. Our study aligns with this trajectory by combining the extraction of features from EEG data with a CNN classifier, with the aim of improving the precision and usability of SSVEP systems.